\section{Word Search}
\label{sec:rec-word-search}

\problemstatement{Given an $m \times n$ grid of characters and a string
$\textit{word}$, determine whether $\textit{word}$ can be formed by a
path of \textbf{sequentially adjacent} cells (horizontally or vertically
neighboring, not diagonal), where the \textbf{same cell may not be used
more than once} within the path.}

\begin{patternbox}
This is Rat in a Maze's (Section~\ref{sec:rec-rat-in-maze}) grid
backtracking template, combined with Sudoku's
(Section~\ref{sec:rec-sudoku-solver}) boolean-return,
short-circuit-on-success control flow: we need a path through the grid
(mark-before-recurse, unmark-after-return, exactly as in the maze), but
we only need to know \emph{whether one exists} matching the word, so the
search should stop the instant a complete match is found, exactly as
Sudoku stops the instant a full board is completed.
\end{patternbox}

\subsection*{Understanding the problem carefully}

Try starting the search from every cell that matches $\textit{word}$'s
first character (since the path could begin anywhere in the grid).
From a starting cell, the path must match $\textit{word}$ character by
character, moving to an adjacent, not-yet-used cell at each step; the
search succeeds the moment all of $\textit{word}$'s characters have been
matched, and fails along a particular path the moment a required
character cannot be matched by any available neighbor.

\begin{ideabox}[Boolean Grid Backtracking, with In-Place Visited Marking]
$\textsc{Solve}(\textit{row}, \textit{col}, \textit{index})$: if
$\textit{index} = |\textit{word}|$: return \texttt{true} (every
character matched). If out of bounds, or
$\textit{grid}[\textit{row}][\textit{col}] \ne \textit{word}[\textit{index}]$,
or this cell is already used in the current path: return
\texttt{false}. Otherwise: temporarily mark this cell as used (a common
space-saving trick: overwrite the grid cell itself with a sentinel
character unlikely to match anything, rather than maintaining a separate
visited grid). Try all four neighboring directions, recursing with
$\textit{index}+1$; if \emph{any} direction returns \texttt{true}: restore
the cell and return \texttt{true} immediately. If none do: restore the
cell and return \texttt{false}.
\end{ideabox}

\subsection*{Why overwriting the cell in place is a safe substitute for a separate visited array}

Since the grid cell's original character is only needed for the
\emph{duration} of exploring paths through it, and we always restore it
before returning control to the caller (whether the exploration
succeeded or failed), no information is permanently lost -- by the time
any sibling branch or ancestor call inspects this cell again, it has
already been restored to its original character. This avoids allocating
an entire second $m \times n$ boolean grid purely to track visited
cells, at the cost of briefly mutating the input grid during the search
(acceptable as long as the grid is restored to its original state by
the time the overall function returns, which the undo step guarantees).

\subsection*{Algorithm}

\begin{enumerate}
  \item For every cell $(\textit{row}, \textit{col})$ with
        $\textit{grid}[\textit{row}][\textit{col}] = \textit{word}[0]$:
        if $\textsc{Solve}(\textit{row}, \textit{col}, 0)$ returns
        \texttt{true}: the word exists; return \texttt{true}.
  \item $\textsc{Solve}(\textit{row}, \textit{col}, \textit{index})$:
  \begin{enumerate}
    \item If $\textit{index} = |\textit{word}|$: return \texttt{true}.
    \item If out of bounds, or character mismatch, or this exact cell
          is the sentinel (already in use): return \texttt{false}.
    \item Save the original character; overwrite the cell with a
          sentinel.
    \item For each of the $4$ directions: if
          $\textsc{Solve}(\textit{newRow}, \textit{newCol},
          \textit{index}+1)$ returns \texttt{true}: restore the cell;
          return \texttt{true}.
    \item Restore the cell; return \texttt{false}.
  \end{enumerate}
  \item If no starting cell succeeds: return \texttt{false}.
\end{enumerate}

\subsection*{Dry run}

\dryrun{Take the grid
$\begin{bmatrix} A&B&C&E \\ S&F&C&S \\ A&D&E&E \end{bmatrix}$,
$\textit{word} = \texttt{"SEE"}$.}

\begin{itemize}
  \item Starting cells matching \texttt{'S'}: $(1,0)$ and $(1,3)$.
  \item Try $(1,0)$: index $0$ matches \texttt{'S'}. Mark used. Try
        neighbors for index $1$ (\texttt{'E'}): up $(0,0)=\texttt{A}$,
        no; down $(2,0)=\texttt{A}$, no; left out of bounds; right
        $(1,1)=\texttt{F}$, no. All fail. Restore $(1,0)$. Return
        \texttt{false}.
  \item Try $(1,3)$: index $0$ matches \texttt{'S'}. Mark used. Try
        neighbors for index $1$ (\texttt{'E'}): up $(0,3)=\texttt{E}$,
        match! Mark used. Try neighbors for index $2$ (\texttt{'E'}):
        up out of bounds; down (back to $(1,3)$, sentinel, fail); left
        $(0,2)=\texttt{C}$, no; right out of bounds. All fail from
        $(0,3)$. Restore $(0,3)$. Try down from $(1,3)$ for index $1$:
        $(2,3)=\texttt{E}$, match! Mark used. Try neighbors for index
        $2$: up (back to $(1,3)$, sentinel, fail); left
        $(2,2)=\texttt{E}$, match! index $2=3=|\textit{word}|$? Index
        becomes $3$ which equals $|\textit{word}|=3$: return
        \texttt{true} immediately.
  \end{itemize}

The word \texttt{"SEE"} is found via the path $(1,3) \to (2,3) \to
(2,2)$.

\subsection*{C++ implementation}

\begin{lstlisting}
#include <bits/stdc++.h>
using namespace std;

bool solve(int row, int col, int index, vector<vector<char>>& grid,
           string& word) {
    if (index == (int)word.size()) return true;

    int m = grid.size(), n = grid[0].size();
    if (row < 0 || row >= m || col < 0 || col >= n) return false;
    if (grid[row][col] != word[index]) return false;

    char original = grid[row][col];
    grid[row][col] = '#'; // sentinel: marks this cell as in use

    int dr[] = {1, -1, 0, 0};
    int dc[] = {0, 0, 1, -1};
    bool found = false;
    for (int d = 0; d < 4 && !found; d++) {
        found = solve(row + dr[d], col + dc[d], index + 1, grid, word);
    }

    grid[row][col] = original; // restore
    return found;
}

bool wordSearch(vector<vector<char>>& grid, string word) {
    int m = grid.size(), n = grid[0].size();
    for (int r = 0; r < m; r++) {
        for (int c = 0; c < n; c++) {
            if (grid[r][c] == word[0] && solve(r, c, 0, grid, word)) {
                return true;
            }
        }
    }
    return false;
}

int main() {
    vector<vector<char>> grid = {
        {'A','B','C','E'},
        {'S','F','C','S'},
        {'A','D','E','E'}
    };
    cout << (wordSearch(grid, "SEE") ? "true" : "false") << endl;
    return 0;
}
\end{lstlisting}

\begin{complexitybox}
\textbf{Time Complexity.} $O(mn \cdot 4^L)$, where $L$ is the word's
length -- $mn$ possible starting cells, each potentially exploring up to
$4$ directions at every one of $L$ steps.
\textbf{Space Complexity.} $O(L)$ for the recursion depth (no separate
visited grid needed, thanks to the in-place sentinel trick).
\end{complexitybox}

\begin{takeawaybox}
When a grid-backtracking problem only needs a single existence answer,
combine Rat in a Maze's mark/unmark mechanics with Sudoku's
short-circuit-on-success boolean propagation -- and consider whether the
``visited'' marking can be folded directly into the grid itself (via a
temporary sentinel overwrite) rather than allocating a separate
tracking structure, whenever the grid's cells are not otherwise needed
during the brief window before they are restored.
\end{takeawaybox}
